\documentclass{article}
\usepackage{lmodern}
\usepackage{amsmath}
\usepackage{listings}
\usepackage{fullpage}
\usepackage{parskip}
\usepackage{xcolor}
\lstset{
	basicstyle=\ttfamily,
	columns=fullflexible,
	frame=single,
	breaklines=true,
	postbreak=\mbox{\textcolor{red}{$\hookrightarrow$}\space},
}
\begin{document}
\title{Experiment `p\_6\_2\_d\_replace\_with\_larger\_neighbor\_seq' Results}
\author{\today}
\date{}
\maketitle
\textbf{Experiment outcome: }SUCCESS
\\\textbf{Bad responses: }0
\\\textbf{Responses containing}~\texttt{assume}~\textbf{: }0
\\\textbf{Responses containing}~\texttt{expect}~\textbf{: }0
\\\textbf{Resolution attempts: }1
\\\textbf{Hard fails (resolution): }0
\\\textbf{Soft fails (resolution): }0
\\\textbf{Verification attempts: }1
\section*{Problem Specification}
\textbf{Problem name: }p\_6\_2\_d\_replace\_with\_larger\_neighbor\_seq
\\\textbf{Natural language statement: }null
\\\textbf{Method signature: }p\_6\_2\_d\_replace\_with\_larger\_neighbor\_seq(inputs: seq<int>) returns (result: seq<int>)
\subsection*{Ensures}
\begin{itemize}
\item \texttt{|result| == |inputs|}
\item \texttt{result[0] == inputs[0]}
\item \texttt{result[|inputs| - 1] == inputs[|inputs| - 1]}
\item \texttt{forall i :: 1 <= i < |inputs| - 1 ==> result[i] == if inputs[i - 1] >= inputs[i + 1] then inputs[i - 1] else inputs[i + 1]}
\end{itemize}
\subsection*{Requires}
\begin{itemize}
\item \texttt{|inputs| >= 3}
\end{itemize}
\clearpage
\section*{GenAI interactions}
Below you will find all interactions between the `user' (program) and the `assistant' (OpenAI).
\subsection*{Program $\rightarrow$ GenAI}
\begin{lstlisting}
You are given the following task to perform in Dafny:

Write a method to satisfy the below requirements.

The signature should be:

p_6_2_d_replace_with_larger_neighbor_seq(inputs: seq<int>) returns (result: seq<int>)

The method should respect the following contract:

requires |inputs| >= 3, ensures |result| == |inputs|, ensures result[0] == inputs[0], ensures result[|inputs| - 1] == inputs[|inputs| - 1], ensures forall i :: 1 <= i < |inputs| - 1 ==> result[i] == if inputs[i - 1] >= inputs[i + 1] then inputs[i - 1] else inputs[i + 1]

Produce and show only the Dafny body of this method, including the curly braces that surround it. Do not show the signature nor contract. You must not use 'assume' nor 'expect' anywhere in your code.
\end{lstlisting}
\subsection*{GenAI $\rightarrow$ Program}
\textbf{System fingerprint: }null
\\\textbf{ID: }chatcmpl-DnNmR3vlHfplvuYkiRsFwBmRm3jfS
\\\textbf{Model: }gpt-5.5-2026-04-23
\\\textbf{Created at: }1780661855
\\\textbf{Finish reason: }stop
\\\textbf{Usage: }434 tokens in, and 735 tokens out
\begin{lstlisting}
{
  var n := |inputs|;
  result := [inputs[0]];
  var i := 1;
  while i < n - 1
    invariant n == |inputs|
    invariant 1 <= i <= n - 1
    invariant |result| == i
    invariant result[0] == inputs[0]
    invariant forall j :: 1 <= j < i ==> result[j] == if inputs[j - 1] >= inputs[j + 1] then inputs[j - 1] else inputs[j + 1]
    decreases n - 1 - i
  {
    var v := if inputs[i - 1] >= inputs[i + 1] then inputs[i - 1] else inputs[i + 1];
    result := result + [v];
    i := i + 1;
  }
  result := result + [inputs[n - 1]];
}
\end{lstlisting}
\section*{Final Program}
\begin{lstlisting}
method p_6_2_d_replace_with_larger_neighbor_seq(inputs: seq<int>) returns (result: seq<int>)
	requires |inputs| >= 3
	ensures |result| == |inputs|
	ensures result[0] == inputs[0]
	ensures result[|inputs| - 1] == inputs[|inputs| - 1]
	ensures forall i :: 1 <= i < |inputs| - 1 ==> result[i] == if inputs[i - 1] >= inputs[i + 1] then inputs[i - 1] else inputs[i + 1]
{
  var n := |inputs|;
  result := [inputs[0]];
  var i := 1;
  while i < n - 1
    invariant n == |inputs|
    invariant 1 <= i <= n - 1
    invariant |result| == i
    invariant result[0] == inputs[0]
    invariant forall j :: 1 <= j < i ==> result[j] == if inputs[j - 1] >= inputs[j + 1] then inputs[j - 1] else inputs[j + 1]
    decreases n - 1 - i
  {
    var v := if inputs[i - 1] >= inputs[i + 1] then inputs[i - 1] else inputs[i + 1];
    result := result + [v];
    i := i + 1;
  }
  result := result + [inputs[n - 1]];
}
\end{lstlisting}
\section*{Total Token Usage}
\textbf{Input tokens: }434
\\\textbf{Output tokens: }735
\\\textbf{Reasoning tokens: }512
\\\textbf{Sum of `total tokens': }1169
\section*{Experiment Timings}
\textbf{Overall Experiment} started at 1780661853698, ended at 1780661893148, lasting 39450ms (39.45 seconds)
\\\textbf{Iteration \#1} started at 1780661853716, ended at 1780661893148, lasting 39432ms (39.43 seconds)
\end{document}
